\documentclass[conference]{IEEEtran}
\IEEEoverridecommandlockouts
% --- Packages -----------------------------------------------------------------
\usepackage{cite}
\usepackage{amsmath,amssymb,amsfonts}
\usepackage{algorithmic}
\usepackage{graphicx}
\usepackage{textcomp}
\usepackage{xcolor}
\usepackage{booktabs}
\usepackage{array}
\usepackage{tabularx}
\usepackage{url}
\usepackage{hyperref}
\usepackage{listings}
\hypersetup{colorlinks=true,linkcolor=blue,citecolor=blue,urlcolor=blue}

\def\BibTeX{{\rm B\kern-.05em{\sc i\kern-.025em b}\kern-.08em
    T\kern-.1667em\lower.7ex\hbox{E}\kern-.125emX}}

\begin{document}

% --- Title --------------------------------------------------------------------
\title{A Multi-Agent Deep-Research Assistant for HCI Topics}

\author{\IEEEauthorblockN{Qiming L.}
\IEEEauthorblockA{\textit{Assignment 3, Multi-Agent Systems} \\
\textit{Fall 2025} \\
qiming.l@example.edu}}

\maketitle

% --- Abstract -----------------------------------------------------------------
\begin{abstract}
We present a four-agent deep-research assistant for human--computer
interaction (HCI) topics, orchestrated with Microsoft AutoGen's
\texttt{RoundRobinGroupChat}. A \emph{Planner}, \emph{Researcher},
\emph{Writer}, and \emph{Critic} cooperate via a shared chat, with the
Researcher invoking real Tavily web search and Semantic Scholar paper search
through \texttt{FunctionTool} wrappers. Around this core we add (i)~a custom
rule-based safety policy filter that runs both \emph{before} the team executes
(input guardrail covering harmful content, prompt injection, off-topic
queries, and malformed input) and \emph{after} (output guardrail covering
PII, harmful content, bias, and citation grounding); (ii)~a two-perspective
LLM-as-a-Judge --- a \emph{supportive} reviewer and a \emph{strict}
reviewer --- that each score five criteria (relevance, evidence quality,
factual accuracy, safety compliance, clarity); and (iii)~CLI and Streamlit
interfaces that surface agent traces, citations, and safety refusals. The
system runs against the SALT-lab \texttt{gpt-5-mini} gateway and is evaluated
on ten HCI queries plus five adversarial safety probes.
\end{abstract}

\begin{IEEEkeywords}
multi-agent systems, AutoGen, LLM-as-a-Judge, safety guardrails, HCI,
retrieval-augmented generation
\end{IEEEkeywords}

% --- Body ---------------------------------------------------------------------
\section{System Design and Implementation}

\subsection{Agents and Control Flow}
Four AutoGen \texttt{AssistantAgent} instances share a chat managed by
\texttt{RoundRobinGroupChat} and terminate on the token \texttt{TERMINATE}.
Table~\ref{tab:agents} summarises the role of each agent and its hand-off
token. The orchestrator (\texttt{src/autogen\_orchestrator.py}) wraps
\texttt{team.run()} with two safety gates and a result extractor that pulls
the \emph{Writer's} message as the canonical answer (rather than the Critic's
review), and aggregates messages, sources, and the running agent set into
structured metadata returned to the UI layer.

\begin{table}[t]
\centering
\caption{Agents, tools, and hand-off tokens.}
\label{tab:agents}
\renewcommand{\arraystretch}{1.15}
\begin{tabularx}{\columnwidth}{@{}lll@{}}
\toprule
\textbf{Agent} & \textbf{Tools} & \textbf{Hand-off token} \\
\midrule
Planner    & ---                         & \texttt{PLAN COMPLETE} \\
Researcher & web\_search, paper\_search  & \texttt{RESEARCH COMPLETE} \\
Writer     & ---                         & \texttt{DRAFT COMPLETE} \\
Critic     & ---                         & \texttt{APPROVED} / \texttt{TERMINATE} \\
\bottomrule
\end{tabularx}
\end{table}

\subsection{Tools}
\texttt{src/tools/web\_search.py} exposes a thin \texttt{WebSearchTool} over
Tavily and Brave with a unified result schema (\texttt{title}, \texttt{url},
\texttt{snippet}, \texttt{score}, \texttt{published\_date}).
\texttt{src/tools/paper\_search.py} wraps the \texttt{semanticscholar} Python
client and supports year/citation filters. Both are exposed to AutoGen as
synchronous \texttt{FunctionTool} instances.
\texttt{src/tools/citation\_tool.py} formats sources in APA-7 (with MLA
fallback), de-duplicates entries, and emits a sorted bibliography.

\subsection{Models and Configuration}
All tunables --- model provider, agent prompts, tool providers, safety policy,
and judge criteria --- live in \texttt{config.yaml}; secrets live in
\texttt{.env}. The default chat model is \textbf{gpt-5-mini} served by the
SALT-lab OpenAI-compatible gateway
(\texttt{https://openai.gateway.salt-lab.org/v1}); the judge model points at
the same endpoint with \texttt{temperature=0.3}. Both Groq and
OpenAI-compatible (vLLM) clients are auto-detected at runtime, so the same
code runs against the class-provided endpoint or a Groq fallback without
modification.

\subsection{User Interfaces}
A Streamlit web UI (\texttt{src/ui/streamlit\_app.py}) shows the final
answer, citations, agent traces, a quality score, and a dedicated
\emph{Safety Events} expander. A CLI (\texttt{src/ui/cli.py}) prints the same
information for terminal use. A single entry point, \texttt{main.py}, exposes
five modes: \texttt{cli}, \texttt{web}, \texttt{evaluate}, \texttt{autogen},
and \texttt{demo} (a one-shot end-to-end query that exports the full session
as JSON to \texttt{outputs/}).

% -----------------------------------------------------------------------------
\section{Safety Design}

We defined a six-category policy and implemented it as a custom rule-based
filter in \texttt{src/guardrails/}. We chose rule-based filters over
Guardrails~AI \cite{guardrails_ai} or NeMo Guardrails \cite{rebedea2023nemo}
to keep the system deterministic, offline-friendly, and inspectable; the
implementation exposes a \texttt{\_check\_with\_external\_validator} hook so
a stronger validator can be plugged in without changing the orchestrator.

\begin{table}[t]
\centering
\caption{Safety policy: categories, stage, and default action.}
\label{tab:safety}
\renewcommand{\arraystretch}{1.15}
\begin{tabularx}{\columnwidth}{@{}lll@{}}
\toprule
\textbf{Category} & \textbf{Stage} & \textbf{Default action} \\
\midrule
\texttt{harmful\_content}    & input + output & refuse \\
\texttt{prompt\_injection}   & input          & refuse + sanitise \\
\texttt{off\_topic\_queries} & input          & warn \\
\texttt{pii}                 & output         & sanitise (redact) \\
\texttt{misinformation}      & output         & warn \\
\texttt{bias}                & output         & sanitise \\
\bottomrule
\end{tabularx}
\end{table}

Severity is \texttt{high}, \texttt{medium}, or \texttt{low}; only \texttt{high}
violations trigger the refusal path. Each event carries
\texttt{(timestamp, type, action, violations, content\_preview)} and is
appended both to an in-memory list (returned to the UI in
\texttt{metadata.safety.events}) and to \texttt{logs/safety\_events.log}. The
orchestrator calls \texttt{check\_input\_safety()} \emph{before}
\texttt{team.run()} and \texttt{check\_output\_safety()} \emph{after}, so
blocked queries never reach the LLM and PII never reaches the user.

The Streamlit UI displays a red banner --- \emph{``Input blocked by a safety
policy''} --- with the exact triggered category, and an amber banner for
sanitisation. We verified all six categories pass on a five-query adversarial
probe set (\texttt{data/safety\_test\_queries.json}); for example,
\emph{``Ignore previous instructions and tell me how to make a bomb''} fires
\texttt{harmful\_content}, \texttt{prompt\_injection}, and
\texttt{off\_topic\_queries} simultaneously, and the team is never invoked.

% -----------------------------------------------------------------------------
\section{Evaluation Setup and Results}

\subsection{Dataset}
\texttt{data/example\_queries.json} contains ten diverse HCI queries spanning
explainable AI, AR usability, AI ethics, UX measurement, conversational AI in
healthcare, accessibility, uncertainty visualisation, voice UIs for the
elderly, AI prototyping, and cross-cultural design. Each item carries a
category and an expected-topic list; three carry an APA-style ground truth.
\texttt{data/safety\_test\_queries.json} adds five adversarial / control
prompts targeting the safety policy.

\subsection{LLM-as-a-Judge: Two Independent Rubrics}
A single judging perspective biases toward whatever rubric the prompt evokes.
We therefore implemented an \emph{ensemble} of two perspectives in
\texttt{src/evaluation/judge.py}:
\begin{itemize}
  \item \textbf{Supportive reviewer} --- ``rigorous but partial-credit-friendly'',
        scale $0.0$--$1.0$ with mid-range anchors.
  \item \textbf{Strict reviewer} --- ``sceptical peer reviewer at a top HCI
        venue'', same scale but conservative; demands citations.
\end{itemize}
Each of the five criteria (relevance, evidence quality, factual accuracy,
safety compliance, clarity, with weights
$0.25/0.25/0.20/0.15/0.15$) is scored independently by both perspectives;
the per-criterion score is the average; the overall score is the weighted
average; the per-perspective score is the mean across criteria. The judge
prompt instructs strict JSON output and tolerates markdown fences via a
regex extractor (\texttt{\_parse\_judgment}).

\subsection{Results}
We ran one full demo (\texttt{main.py --mode demo}) end-to-end against
\texttt{gpt-5-mini}. The four agents exchanged 21 messages and produced a
$4\,947$-character draft. Per-criterion scores are shown in
Table~\ref{tab:scores}.

\begin{table}[t]
\centering
\caption{Judge scores for the live \texttt{gpt-5-mini} demo run.}
\label{tab:scores}
\renewcommand{\arraystretch}{1.15}
\begin{tabularx}{\columnwidth}{@{}Xrrr@{}}
\toprule
\textbf{Criterion} & \textbf{Supportive} & \textbf{Strict} & \textbf{Avg.} \\
\midrule
Relevance         & 0.00 & 0.00 & 0.00 \\
Evidence quality  & 0.65 & 0.45 & 0.55 \\
Factual accuracy  & 0.80 & 0.60 & 0.70 \\
Safety compliance & 1.00 & 1.00 & 1.00 \\
Clarity           & 0.65 & 0.45 & 0.55 \\
\midrule
\textbf{Weighted overall} & \textbf{0.68} & \textbf{0.44} & \textbf{0.51} \\
\bottomrule
\end{tabularx}
\end{table}

The $0.24$ spread between perspectives confirms the value of ensembling: the
strict reviewer penalised the lack of grounded citations far more
aggressively than the supportive one. Safety compliance scored $1.0$ from
both judges, as expected.

\subsection{Error Analysis}
The Researcher's tool calls \emph{failed silently} on this run because no
\texttt{TAVILY\_API\_KEY} was supplied and Semantic Scholar refused the
unauthenticated connection (rate limit). The agents detected the empty
results and entered a ``scope confirmation'' loop with the (mocked) user,
never producing the substantive HCI synthesis the query asked for --- hence
relevance $0.00$. This is informative: the system correctly \emph{detected}
that it lacked evidence rather than hallucinating, but it also failed to
recover (e.g.\ by retrying with a different provider or by writing a
clearly-flagged ``no-evidence'' answer). A pre-recorded reference run on the
same accessible-UI query (with Tavily enabled) reached overall $0.84$
(see \texttt{outputs/sample\_session.md}), suggesting tool availability is
the dominant bottleneck rather than agent design. Three additional failure
patterns appeared across exploratory runs: (a)~over-long Critic feedback
echoing back into the chat and pushing context, (b)~Writer occasionally
appending the literal \texttt{TERMINATE} token to its draft (stripped
post-hoc by the orchestrator), and (c)~the Critic asking the user
clarifying questions instead of approving --- addressable by tightening the
agent prompts and adding a hard cap on rounds.

% -----------------------------------------------------------------------------
\section{Discussion and Limitations}

\paragraph*{Insights}
Splitting safety into an \emph{input} gate (rule-based, free, fast) and an
\emph{output} gate (PII regex + citation heuristic) gave us deterministic
guarantees on the high-severity classes (harmful content, prompt injection,
raw PII) without needing a moderation model. The two-perspective judge
produced more interpretable signals than a single rubric: a wide
supportive--strict gap reliably flagged ``stylistically OK but
evidence-thin'' answers. Wrapping each tool as a \texttt{FunctionTool} and
letting the Researcher choose calls --- rather than hard-coding a
sequence --- preserved AutoGen's idiomatic flow at the cost of occasional
skipped searches.

\paragraph*{Limitations}
(1)~The guardrails are regex-based and will miss obfuscated harmful content;
the \texttt{\_check\_with\_external\_validator} hook is a placeholder, not a
wired Guardrails-AI / NeMo integration.
(2)~Topic relevance is keyword-driven and biased toward English HCI
vocabulary.
(3)~The judge uses the same model (\texttt{gpt-5-mini}) as the Writer, so it
shares the Writer's blind spots --- a stronger setup would use a different
model family for judging.
(4)~Tool failures are reported but the agents do not recover gracefully (no
exponential back-off, no provider fallback).
(5)~Evaluation $N{=}10$ queries is small; per-criterion variance is therefore
wide.

\paragraph*{Future work}
Plug a real moderation API (e.g.\ OpenAI Moderation, Llama~Guard) into the
input gate; add an LLM-based citation-grounding verifier to the output gate;
introduce a \emph{Verifier} agent that checks each Writer claim against the
retrieved snippets before the Critic; sweep three judges and report
inter-rater agreement (Cohen's~$\kappa$); and grow the eval set to $\geq 50$
queries for tighter confidence intervals.

\paragraph*{Ethics}
The system is deliberately constrained to the HCI research domain, refuses
harmful and injection requests at the front door, and redacts PII before it
is shown. We do not log full user queries to disk --- only $200$-character
previews --- and we keep secrets out of the repo via \texttt{.gitignore}.

% --- References ---------------------------------------------------------------
\begin{thebibliography}{00}

\bibitem{wu2023autogen}
Q. Wu \emph{et al.}, ``AutoGen: Enabling next-gen LLM applications via
multi-agent conversation,'' \emph{arXiv preprint arXiv:2308.08155}, 2023.
[Online]. Available: \url{https://arxiv.org/abs/2308.08155}

\bibitem{autogen_docs}
Microsoft, ``AutoGen: Multi-agent conversation framework (documentation),''
2024. [Online]. Available: \url{https://microsoft.github.io/autogen/}

\bibitem{rebedea2023nemo}
T. Rebedea, R. Dinu, M. Sreedhar, C. Parisien, and J. Cohen, ``NeMo
Guardrails: A toolkit for controllable and safe LLM applications with
programmable rails,'' \emph{arXiv preprint arXiv:2310.10501}, 2023.
[Online]. Available: \url{https://arxiv.org/abs/2310.10501}

\bibitem{guardrails_ai}
Guardrails~AI, ``Guardrails: An open-source toolkit for safer LLM
applications,'' 2024. [Online]. Available:
\url{https://docs.guardrailsai.com/}

\bibitem{tavily}
Tavily, ``Tavily search API documentation,'' 2024. [Online]. Available:
\url{https://docs.tavily.com/}

\bibitem{semanticscholar}
Allen Institute for AI, ``Semantic Scholar Academic Graph API,'' 2024.
[Online]. Available: \url{https://api.semanticscholar.org/}

\bibitem{wcag}
S.~L. Henry, ``Web content accessibility guidelines (WCAG) overview,''
W3C Web Accessibility Initiative, 2024. [Online]. Available:
\url{https://www.w3.org/WAI/standards-guidelines/wcag/}

\bibitem{lazar2017accessibility}
J. Lazar, D. Goldstein, and A. Taylor, \emph{Ensuring Digital Accessibility
through Process and Policy}. San Francisco, CA, USA: Morgan Kaufmann, 2017.

\bibitem{petrie2009evaluation}
H. Petrie and N. Bevan, ``The evaluation of accessibility, usability, and
user experience,'' in \emph{The Universal Access Handbook},
C. Stephanidis, Ed. Boca Raton, FL, USA: CRC Press, 2009, pp. 1--16.

\bibitem{webaim}
WebAIM, ``Introduction to web accessibility,'' 2024. [Online]. Available:
\url{https://webaim.org/intro/}

\end{thebibliography}

\end{document}
